\section{Introduction}

Urban air pollution remains one of the most critical public health and environmental challenges of the twenty-first century.
The World Health Organization estimates that approximately 99\% of the global population
breathes air containing pollutant concentrations above guideline limits, and the
\textit{State of Global Air 2025} report attributes 7.9 million deaths in 2023 to air pollution
worldwide \cite{world_health_organization_air_nodate,health_effects_institute_state_2025}.
Transportation is an important contributor to this burden. The IPCC Sixth Assessment Report
states that the transport sector accounted for 23\% of global energy-related
$\mathrm{CO_2}$ emissions in 2019, with road vehicles responsible for the majority of that total
\cite{ipcc_climate_2022,creutzig_transport_2015}. In urban environments, these emissions are
especially consequential because they are released close to where people live, work, and travel
\cite{uherek_transport_2010}.

New York City illustrates this problem clearly. It combines a dense population, heavy roadway
activity, and substantial exposure to traffic-related air pollution
\cite{nyc_department_of_city_planning_population_2024,
nyc_department_of_health_and_mental_hygiene_traffic_nodate}.
According to the NYC Department of Health and Mental Hygiene, fine particulate matter 
(PM2.5) is the most harmful air pollutant for humans to breathe, and everyday car, bus, 
and truck traffic accounts for 14\% of PM2.5 emitted from NYC activities
\cite{nyc_department_of_health_and_mental_hygiene_traffic_nodate}. At the same time,
vehicle emissions are not distributed uniformly across the urban environment. They vary by
location, time of day, congestion, and, importantly, by vehicle type. Heavy vehicles, for example,
contribute disproportionately to pollutants such as nitrogen oxides and particulate matter relative
to light-duty passenger vehicles \cite{wang_high-resolution_2025}. As a result, estimating
localized transportation emissions not only requires the total traffic volume, but also 
information about the composition of traffic flow.

Obtaining this type of information continuously and at fine spatial scales remains difficult.
Traditional traffic-monitoring methods, such as manual counts and Automated Traffic Recorders
(ATRs), remain essential for transportation planning, but they are often collected over limited
durations and at limited locations \cite{FHWA_traffic_2010,
NYC_open_data_atr}. Likewise, fixed air-quality monitoring stations measure
pollutant concentrations in the air, but they do not directly identify which vehicle classes are present
at a given place and time. These limitations matter because localized emissions burdens are shaped
not only by how many vehicles are present, but by what kinds of vehicles are moving through a
corridor and under what traffic conditions.

Existing traffic-camera infrastructure offers one possible way to address this gap. New York City
and New York State maintain large networks of public traffic cameras, and the 511NY system
provides access to many live roadway images without requiring new sensing hardware
\cite{NYSDOT_c038145_2025, NYCDOT_realtime, 511ny_api_docs}.
These cameras were originally deployed for traffic operations and incident monitoring,
and the imagery they provide is challenging for computer vision. The images from the NYC CCTV traffic cameras
used in this study have a resolution of 352$\times$240 pixels and often contain distant vehicles, motion blur,
occlusion, glare, rain, and non-ideal viewing angles. Even with the low-image quality, 
this infrastructure offers an appealing foundation for scalable observation if 
vehicles can be detected and categorized reliably enough for downstream analysis.

To examine this problem, a vehicle-detection pipeline was developed using publicly available
New York City traffic-camera imagery and evaluated for emissions-oriented traffic measurement.
The approach is built around a YOLOv11-based detector trained on a custom annotated dataset of
low-resolution urban traffic scenes. Rather than using a generic vehicle taxonomy, the detector is
trained and evaluated with an EPA MOtor Vehicle Emission Simulator (MOVES) based class 
structure designed to balance what can be distinguished reliably in low-resolution imagery
with what remains useful for downstream emissions analysis 
\cite{us_epa_moves3_2020}. Ambiguous vehicles are handled explicitly
through an \verb|unknown_vehicle| label during annotation and through ignore-aware training 
and evaluation procedures, rather than being forced into potentially incorrect supervised classes.

The experimental results show that the final YOLOv11s detector outperformed the Faster R-CNN
and DETR-style alternatives evaluated in this study, and that the most meaningful gains came from
targeted data-centric refinement rather than broader tuning alone. The final model was then assessed
not only on a held-out test set, but also on unseen traffic cameras and through comparison with ATR
hourly traffic counts. These evaluations were included because the goal of the work was not simply
to maximize object-detection benchmark performance, but to determine whether traffic-camera
imagery can support practical count-based workflows relevant to emissions estimation.
The results suggest that low-resolution public traffic-camera imagery, despite its
visual limitations, can still support useful vehicle counting and categorization when the detection
task is formulated around the realities of the imagery and the needs of the downstream application.
In this sense, the proposed pipeline is positioned not as a replacement for existing traffic or
air-quality monitoring systems, but as a complementary framework for producing more spatially
and temporally resolved information about traffic composition from infrastructure that already exists.
